% Lösungen Orthogonalität (zusammenbau v0.4, Heft)
\einheitenkopf{Orthogonalität}

\erg{81}{a) $\overrightarrow{BA}$ und $\overrightarrow{BC}$ – beide Schenkel gehen vom Scheitel $B$ aus \quad b) $\overrightarrow{AB} = (3 | 2 | 0)$, $\overrightarrow{AC} = (-2 | 3 | 0)$; $\overrightarrow{AB} \circ \overrightarrow{AC} = 3 \cdot (-2) + 2 \cdot 3 + 0 \cdot 0 = 0$ – rechter Winkel bei $A$ \quad c) $OP$ liegt auf der $x$-Achse, $OQ$ in der $yz$-Ebene; die $x$-Achse steht senkrecht auf der $yz$-Ebene – rechter Winkel bei $O$}
\erg{82}{a) $\overrightarrow{OA} \circ \overrightarrow{OB} = 0 + 0 + 0 = 0$; $|\overrightarrow{OB}| = 13$, $|\overrightarrow{OA}| = \sqrt{25 + 144} = 13$ \quad b) $\overrightarrow{CB} \circ \overrightarrow{CE} = 0 \cdot (-9) + (-6) \cdot 0 + 0 \cdot 4 = 0$ – rechter Winkel bei $C$; Oberfläche: Rechteck $54$, Dreiecke $CDE$ $18$ und $ADE$ $12$ (rechtwinklig bei $D$), $BCE$ $\frac{1}{2} \cdot 6 \cdot \sqrt{97}$ und $ABE$ $\frac{1}{2} \cdot 9 \cdot \sqrt{52}$ (rechtwinklig bei $A$); zusammen $\approx 146$ cm² \quad c) $\overrightarrow{CD} \circ \overrightarrow{CE} = (-7) \cdot 0 + 0 \cdot (-7) + 0 \cdot 5 = 0$; Katheten $|CD| = 7$ und $|CE| = \sqrt{74}$; Oberfläche: Quadrat $49$, $ABE$ und $BCE$ je $17{,}5$, $CDE$ und $ADE$ je $\frac{1}{2} \cdot 7 \cdot \sqrt{74}$; zusammen $\approx 144{,}2$ cm² \quad d) $\overrightarrow{AB} \circ \overrightarrow{AC_t} = 4 \cdot 3t + 3 \cdot (-4t) + 5 \cdot 0 = 12t - 12t = 0$ für jedes $t$ \quad e) $\overrightarrow{BA} = (-3 | 0 | 2)$, $\overrightarrow{BC_t} = (2t | 0 | 3t)$; $\overrightarrow{BA} \circ \overrightarrow{BC_t} = -6t + 0 + 6t = 0$ für jedes $t$ \quad f) $\overrightarrow{BA} \circ \overrightarrow{BC} = (-\sqrt{7}) \cdot (-\sqrt{7}) + (-6) \cdot 6 + 0 = 7 - 36 = -29 \ne 0$ – bei $B$ kein rechter Winkel \quad g) $\overrightarrow{BA} \circ \overrightarrow{BC} = (-\sqrt{5}) \cdot (-\sqrt{5}) + (-2) \cdot 3 + 0 = 5 - 6 = -1 \ne 0$ – bei $B$ kein rechter Winkel}
\erg{83}{a) null setzen – ein Wert ist gesucht, Gleichung lösen \quad b) $\overrightarrow{AB} \circ \overrightarrow{AC} = (3 | 1 | 0) \circ (1 | t - 2 | 2) = 3 + t - 2 = t + 1 = 0$, also $t = -1$ \quad c) $\overrightarrow{AG} \circ \overrightarrow{CE} = (-12 | 9 | h) \circ (12 | -9 | h) = -144 - 81 + h^2 = 0$, also $h = 15$ ($h = -15$ entfällt); $V = 12 \cdot 9 \cdot 15 = 1620$; $O = 2 \cdot (108 + 180 + 135) = 846$}
\erg{84}{a) $\overrightarrow{BO} \circ \overrightarrow{BA} = (-2 | -2 | -3) \circ (2 | -2 | a - 3) = -4 + 4 - 3a + 9 = 9 - 3a = 0$, also $a = 3$ \quad b) $\overrightarrow{Q_tO} \circ \overrightarrow{Q_tP} = (-2 | -t | -8) \circ (2 | -t | 4) = -4 + t^2 - 32 = t^2 - 36 = 0$, also $t = 6$ oder $t = -6$ \quad c) $\overrightarrow{S_tF} \circ \overrightarrow{S_tM} = (-t | 0 | 2) \circ (5 - t | 1 | 2) = -5t + t^2 + 4 = 0$, also $t = 1$ oder $t = 4$ \quad d) $C(x | y | 5)$: $\overrightarrow{AB} \circ \overrightarrow{AC} = 3(x - 1) + 2(y - 2) + 2 \cdot 5 = 3x + 2y +3 = 0$ und $2x + y = 0$; Lösung $x = 3$, $y = -6$, also $C(3 | -6 | 5)$ \quad e) $C(x | y | -1)$: $\overrightarrow{AB} \circ \overrightarrow{AC} = 1(x - 0) + 3(y - 1) + -1 \cdot -3 = 1x + 3y +0 = 0$ und $x + y = 2$; Lösung $x = 3$, $y = -1$, also $C(3 | -1 | -1)$}
\erg{85}{a) Gerade gegen Ebene – Richtungsvektor kollinear zum Normalenvektor \quad b) $(1 | 1 | 1) \circ (2 | -t | 3) = 2 - t + 3 = 0$, also $t = 5$}
\erg{86}{a) $\overrightarrow{BE} = (4 | -2 | 4) = 2 \cdot (2 | -1 | 2)$, ein Vielfaches des Normalenvektors von $L$ – die Seitenkanten stehen senkrecht auf der Grundfläche, das Prisma ist gerade \quad b) $(2 | 1 | 2) \circ (1 | -4 | 1) = 2 - 4 + 2 = 0$ und $(1 | 0 | -1) \circ (1 | -4 | 1) = 1 + 0 - 1 = 0$ – $\vec n$ steht auf beiden Spannvektoren senkrecht, also auf $E$ \quad c) $\overrightarrow{CG} \circ \vec n = (-3 | 0 | 5) \circ (5 | 5 | 3) = -15 + 0 + 15 = 0$, $\overrightarrow{CF} \circ \vec n = (0 | -3 | 5) \circ (5 | 5 | 3) = 0 - 15 + 15 = 0$; $N\colon 5x + 5y + 3z = 45$ ($B$ eingesetzt ergibt $45$) \quad d) $(0 | 2 | 1) \circ (1 | -1 | 2) = 0 - 2 + 2 = 0$ und $(1 | 1 | 0) \circ (1 | -1 | 2) = 1 - 1 + 0 = 0$ – beide Skalarprodukte null, $\vec n$ ist Normalenvektor \quad e) Richtungsvektor von $h$: $(0 | 1 | 0)$; $(-3 | 0 | 2) \circ (0 | 1 | 0) = 0 + 0 + 0 = 0$ – ja, senkrecht \quad f) Richtungsvektor von $h$: $(1 | 0 | 0)$; $(2 | -3 | 0) \circ (1 | 0 | 0) = 2 \ne 0$ – nicht senkrecht \quad g) $(4 | -3 | 1) \circ (1 | 2 | k) = 4 - 6 + k = 0$, also $k = 2$ \quad h) $(0 | 2 | -5) \circ (3 | a | 2) = 0 + 2a - 10 = 0$, also $a = 5$ \quad i) z. B. $H\colon 3x - 5y - 7z = 0$ – Normalenvektor von $H$ ist der Richtungsvektor $(3 | -5 | -7)$ der Schnittgeraden (er steht auf beiden Normalenvektoren senkrecht), der Stützpunkt ist frei \quad j) z. B. $H\colon 2x + 7y + 3z = 0$ – Normalenvektor ist der Richtungsvektor $(2 | 7 | 3)$ der Schnittgeraden; jede Ebene mit diesem Normalenvektor steht senkrecht auf $E$ und $F$}
\erg{87}{a) $M(2 | 0 | 3)$ ist der Mittelpunkt von $RS$, die Fläche ist $\frac{1}{2} \cdot |RS| \cdot |MT_\lambda|$ mit festem $|RS|$; minimal wird sie bei kleinstem $|MT_\lambda|$, also am Lot von $M$ auf $t$: $\overrightarrow{MT_\lambda} = (\lambda - 2 | 1 + 2\lambda | 0)$ senkrecht zu $(1 | 2 | 0)$; daraus $5\lambda = 0$, $\lambda = 0$}
\erg{88}{a) Mit $M(0 | 3 | 3)$ als Mittelpunkt von $RS$ ist der Flächeninhalt $\frac{1}{2} \cdot |RS| \cdot |MT_\lambda|$; $|RS|$ ist fest, also ist die Fläche am kleinsten, wenn $|MT_\lambda|$ am kleinsten ist – das ist das Lot von $M$ auf $t$, und dort steht $\overrightarrow{MT_\lambda} = (1 + 2\lambda | 0 | \lambda - 3)$ senkrecht auf dem Richtungsvektor $(2 | 0 | 1)$; Lösung $5\lambda - 1 = 0$, $\lambda = 0{,}2$}
